\documentclass[11pt]{article}
\usepackage[margin=1in]{geometry}
\usepackage{booktabs,hyperref}
\title{Bounded Trace-Event Retrieval with Opt-In Payload Preview: A Synthetic Fleet UI Acceptance Study}
\author{AssistX Fleet Research --- Working Draft}
\date{October 9, 2026}
\begin{document}
\maketitle
\begin{abstract}
A progressive timeline frontend can reduce displayed event count while leaving
database read volume, network payload and private field transmission unbounded.
We prototype server-owned, disabled-by-default keyset event metadata pages
bounded to 100 rows and a separate authenticated operator-initiated request for
a 4096-character payload preview. Keyset navigation uses a versioned cursor
over timestamp and event ID; it is not an authenticity token. Synthetic tests
include 1001 events at equal timestamps, cursor corruption, source-group
ownership, and a missing production authentication hook. The addition and
its surrounding contracts passed 131 Python regression tests, while 25 Node
workbench contracts passed. A disconnected Neo4j 5.26.26 instance successfully
compiled and executed both fixed read queries on 1001 generated events. The
legacy unbounded endpoint and existing frontend remain unchanged; neither
production readiness nor source archival completeness is claimed.
\end{abstract}
\section{Motivation and prospectus}
A former workbench render cap of 80 events limited DOM pressure but did not
constrain the full-detail \texttt{get\_trace()} graph read or transport of
\texttt{payload\_json}. The preregistered \texttt{TRACE\_DETAIL\_PAGING\_PROSPECTUS\_20261009.md}
hypothesized keyset pagination, default metadata minimization, source-scoped
opt-in preview, and explicit denial until production authentication and query
admission are established. Neo4j Cypher supports fixed predicates and ordered
projections \cite{cypher}.
\section{Method}
The keyset condition is strictly lexicographic on descending $(ts_{\mathrm{ms}},
event\_id)$: the next page selects events with smaller timestamp or smaller
event ID at equal timestamp. A fixed \texttt{LIMIT \$take} of page size plus
one detects another page; each response includes at most 100 metadata events,
an optional opaque cursor, and \texttt{total\_indexed\_events=null}. Scalar
producer labels and context IDs remain explicitly unverified as physical node
or agent provenance. Payload preview uses a distinct POST with event ID in
the request body, protected by the same existing authentication dependency,
a default-disabled feature flag, and a bounded Cypher substring. Both routes
use a four-second read-only query timeout and \texttt{Cache-Control: no-store}.
\section{Evidence}
\begin{table}[h]\centering
\caption{Synthetic implementation acceptance.}
\begin{tabular}{ll}\toprule
Test & Observation\\\midrule
Python paging and related regression tests & 131 passed\\
Existing JavaScript workbench contract tests & 25 passed\\
1001 equal-timestamp synthetic events & Gap-free fake-session paging\\
Isolated Neo4j 5.26.26, 1001 generated events & Fixed Cypher read queries executed\\
Missing auth injection/flag disabled & HTTP 503, no graph read\\
Live production latency and concurrency & Not measured\\
Frontend switched to bounded server API & Not implemented\\
\bottomrule\end{tabular}\end{table}
The isolated command-line smoke took approximately 2.1--2.3 seconds per
read including Docker and JVM startup; these timings do not establish server
driver latency or production p95. The disconnected container and anonymous
volumes were removed.
\section{Threats to validity and release constraints}
The old unbounded endpoint remains in place and the new routes are not
called by the workbench, so this study does \emph{not} show that current
production transport is capped. Stable-source pagination excludes concurrent
backfills, duplicate event IDs and changes to timestamps. Four-second query
timeouts do not physically fence stuck queries; additional cross-worker
admission and cancellation work remains outstanding. Payload preview has
not undergone independent privacy/redaction approval. The returned metadata
is investigative index evidence, distinct from full-fidelity nontrimming
agent tool-call custody \cite{nist92}. No new operational authority, NAS
mutation, provider routing or model use is permitted by this draft.
\begin{thebibliography}{9}
\bibitem{cypher} Neo4j, \emph{Cypher Manual}, ordering and filtering.
\url{https://neo4j.com/docs/cypher-manual/current/}.
\bibitem{nist92} K. Kent and M. Souppaya, \emph{Guide to Computer Security
Log Management}, NIST SP 800-92 (2006).
\url{https://doi.org/10.6028/NIST.SP.800-92}.
\end{thebibliography}
\end{document}
